\documentclass[crop,tikz,11pt]{standalone}
\usepackage{geometricfigures}
\usepackage{amsmath,amssymb}
\usetikzlibrary{arrows.meta}

\definecolor{net}{HTML}{4E9A4E}

% The three ways of representing M in a single chart. Each panel draws, in its latent plane, the images of the circles
% p_theta = s of M (thin; integer s in (b) and (c), half-integer steps in (a)), the zero section nu = {p_theta = 0} (dotted, as in
% cylinder-to-poloidal), and the hole (open circle, as the apex there).
%   (b) with a cut-off, the round chart of cylinder-to-poloidal, a = 3.2 as there:
%       r^2/2 = a - s, s = -2, ..., 3; the circles reach the apex at s = a, {p_theta >= a} is cut off.
%   (c) symplectic and Euclidean, a symplectic decoder of the whole plane that parametrizes {p_theta < -1}: the chart
%       r^2/2 = -s of rotating-cap (there with the boundary at -1.05), s = -5, ..., -1; the
%       disk r < sqrt(2) inside the circle s = -1 decodes onto the cap off M (net!16, as there).
%   (a) without symplecticity, the homeomorphism r = e^{-s}, s = -1, -1/2, ..., 2 and 3; the circles accumulate at the hole,
%       the strips between them shrink towards it.
% The panels go from no symplecticity to symplectic and Euclidean. All three at the same scale; the strips between circles
% one unit of p_theta apart have area 2pi in (b) and (c), and not in (a). Labels of the hole and of the cap's edge sit below right, on a leader.
\newcommand{\un}{0.5}
\tikzset{lev/.style={fg!45!bg,line width=0.4pt},
         nu/.style={fg!70!bg,densely dotted,line width=0.6pt},
         lab/.style={font=\footnotesize,inner sep=1pt},
         hole/.style={fill=bg,draw=fg,line width=0.45pt}}

\begin{document}
\begin{tikzpicture}[x=\un cm,y=\un cm]

% ---------------------------------------------------------------- (b) with a cut-off: the round chart
\begin{scope}[shift={(9.5,0)}]
  \foreach \s in {-2,-1,1,2,3} {\draw[lev] (0,0) circle ({sqrt(2*(3.2-\s))});}
  \draw[nu] (0,0) circle ({sqrt(6.4)});
  \filldraw[hole] (0,0) circle (1.3pt);
  \node[lab,anchor=south] at (0,3.7) {$p_\theta=-2$};
  \node[lab,anchor=south] at (-0.45,{sqrt(6.4)+0.18}) {$\nu$};
  \draw[fg!50!bg,line width=0.3pt] (0.12,-0.12) -- (2.6,-3.2) node[lab,anchor=north west,fg] {$p_\theta\to a$};
  \node[font=\small] at (0,-4.2) {\textbf{(b)}};
\end{scope}

% ---------------------------------------------------------------- (c) symplectic and Euclidean: whole plane
\begin{scope}[shift={(18.5,0)}]
  \fill[net!16!bg] (0,0) circle ({sqrt(2)});
  \foreach \s in {-5,-4,-3,-2} {\draw[lev] (0,0) circle ({sqrt(-2*\s)});}
  \draw[lev] (0,0) circle ({sqrt(2)});
  \node[lab,anchor=west] at ({sqrt(10)+0.35},0) {$p_\theta=-5$};
  \draw[fg!50!bg,line width=0.3pt] ({sqrt(2)*cos(-50)},{sqrt(2)*sin(-50)}) -- (2.6,-3.2) node[lab,anchor=north west,fg] {$p_\theta=-1$};
  \node[lab,net!70!fg] at (0,0) {cap};
  \node[font=\small] at (0,-4.2) {\textbf{(c)}};
\end{scope}

% ---------------------------------------------------------------- (a) without symplecticity: all of M
\begin{scope}
  \foreach \s in {-1,-0.5,0.5,1,1.5,2,3} {\draw[lev] (0,0) circle ({exp(-\s)});}
  \draw[nu] (0,0) circle (1);
  \filldraw[hole] (0,0) circle (1.3pt);
  \node[lab,anchor=west] at ({exp(1)+0.35},0) {$p_\theta=-1$};
  \node[lab,anchor=south] at (0,1) {$\nu$};
  \draw[fg!50!bg,line width=0.3pt] (0.12,-0.12) -- (2.6,-3.2) node[lab,anchor=north west,fg] {$p_\theta\to\infty$};
  \node[font=\small] at (0,-4.2) {\textbf{(a)}};
\end{scope}

\end{tikzpicture}
\end{document}
